\documentclass{article}

\usepackage[letterpaper,top=2cm,bottom=2cm,left=2.54cm,right=2.54cm,marginparwidth=1.75cm]{geometry}
\usepackage{amsmath}
\usepackage{amssymb}
\usepackage[colorlinks=true, allcolors=blue]{hyperref}
\usepackage{comment}
\usepackage[natbibapa]{apacite}
\usepackage{xcolor} % added for reviewer comments
\usepackage{xr-hyper} % resolve cross-references to Appendix A (eq:lds_state_eq, eq:lyapunov_eq0), clickable
% 1st arg: path to the .aux, relative to this source dir (appendices/), for reading labels.
% 3rd arg: path to the target PDF, relative to THIS pdf's own location (appendices/build/),
% since both compiled PDFs live in the same build/ folder -- hence no "build/" prefix here.
\externaldocument{../build/appendices/appendix_a_lyapunov_derivation/appendix_a_lyapunov_derivation}[appendix_a_lyapunov_derivation.pdf]

\setlength{\parindent}{0pt}

%\title{Appendix B: Analytic Relationship between Patterns of Correlation and Causation for an $N \times N$ Network}

%\author{Yaohui Ding}

\begin{document}

%\maketitle

\section*{Appendix B: Analytic Relationship between Patterns of Correlation and Causation for an $N \times N$ network}
\phantomsection\label{app:appB}
\renewcommand{\theequation}{B-\arabic{equation}}
\setcounter{equation}{0}

\subsection*{The Covariance of Observation $\mathbf{Y} (t)$}

The observation equation of a linear stochastic dynamical system (\ref{eq:lds_state_eq}) is often given by
% seperator %%%%%%%%%%%%%%%%%%%%%%%%%%%%%%%%%%%%%%%%%%%%%%%%%
% H maps the (possibly unobserved) state X(t) to what is actually measured,
% Y(t), corrupted by observation noise V(t).
\begin{equation}
\mathbf{Y}(t)=\mathbf{H}\mathbf{X}(t)+\mathbf{V}(t)
\label{eq:lds_obs_eq}
\end{equation}

where $\mathbf{Y}(t) \in \mathbb{R}^m, \mathbf{V}(t) \in \mathbb{R}^m$, and $\mathbf{H} \in \mathbb{R}^{m \times n}$. Let us assume $n=m$. Additionally,
% seperator %%%%%%%%%%%%%%%%%%%%%%%%%%%%%%%%%%%%%%%%%%%%%%%%%
\begin{equation}
\begin{aligned}
E[\mathbf{V}(t)]&=\mathbf{0},
Cov \left ( \mathbf{V}(t) \right )=\mathbf{\Sigma}_\mathbf{V}
\end{aligned}
\end{equation}

We further assume that the state $\mathbf{X}(t)$ and the observation noise $\mathbf{V}(t)$ are uncorrelated,
% seperator %%%%%%%%%%%%%%%%%%%%%%%%%%%%%%%%%%%%%%%%%%%%%%%%%
\begin{equation}
Cov \left ( \mathbf{X}(t), \mathbf{V}(t) \right ) = E \left[ \mathbf{X}(t) \mathbf{V}(t)^T \right] = \mathbf{0}
\label{eq:state_noise_uncorr}
\end{equation}

Without introducing any ambiguity and to simplify the notation, we will let $\mathbf{X}(t)=\mathbf{X}$ and $\mathbf{V}(t)=\mathbf{V}$. Furthermore, at steady state, $E[\mathbf{X}]=\mathbf{0}$, as established in (\ref{eq:ss_mean_steady}). The mean of $\mathbf{Y}(t)$ at steady state is
% seperator %%%%%%%%%%%%%%%%%%%%%%%%%%%%%%%%%%%%%%%%%%%%%%%
\begin{equation}
\begin{aligned}
E[\mathbf{Y}(t)] &=E \left[ \mathbf{H}\mathbf{X}+\mathbf{V} \right] =E \left[ \mathbf{H}\mathbf{X} \right ] + E \left [ \mathbf{V} \right]
=\mathbf{H} E \left[ \mathbf{X} \right ] + E \left [ \mathbf{V} \right]
=\mathbf{H} * \mathbf{0}+ \mathbf{0}=\mathbf{0}
\end{aligned}
\label{eq:mean_yt}
\end{equation}


% seperator %%%%%%%%%%%%%%%%%%%%%%%%%%%%%%%%%%%%%%%%%%%%%%%
% Expand Y(t)Y(t)^T; the two cross terms vanish by eq:state_noise_uncorr,
% leaving Sigma_Y = H*Sigma_X(infty)*H^T + Sigma_V.
The covariance of $\mathbf{Y}(t)$ is
\begin{equation}
\begin{aligned}
\mathbf{\Sigma_Y} &= E \left[ \left( \mathbf{Y}(t)-E[\mathbf{Y}(t)] \right) \left( \mathbf{Y}(t)-E[\mathbf{Y}(t)] \right)^T \right] \\
&=E \left[ \mathbf{Y}(t) \mathbf{Y}(t)^T \right] \\
&=E \left[ \left( \mathbf{H}\mathbf{X}+\mathbf{V} \right) \left( \mathbf{H}\mathbf{X}+\mathbf{V} \right)^T \right] \\
&=E \left[ \mathbf{H} \mathbf{X} \mathbf{X}^T \mathbf{H}^T + \mathbf{H} \mathbf{X} \mathbf{V}^T + \mathbf{V} \mathbf{X}^T \mathbf{H}^T +\mathbf{V} \mathbf{V}^T \right] \\
&= E \left[ \mathbf{H} \mathbf{X} \mathbf{X}^T \mathbf{H}^T \right] + E \left[ \mathbf{H} \mathbf{X} \mathbf{V}^T \right] + E \left [ \mathbf{V} \mathbf{X}^T \mathbf{H}^T \right] + E \left [\mathbf{V} \mathbf{V}^T \right] \\
&= \mathbf{H} E \left [\mathbf{X} \mathbf{X}^T \right ] \mathbf{H}^T + \mathbf{0} + \mathbf{0} + \mathbf{\Sigma}_V \\
&= \mathbf{H} \mathbf{\Sigma}_{\mathbf{X}(\infty)} \mathbf{H}^T + \mathbf{\Sigma}_V
\end{aligned}
\label{eq:cov_yt}
\end{equation}

\subsection*{Solving the Lyapunov Equation in Appendix A}
% Strategy: apply vec() to turn the matrix equation A*Sigma+Sigma*A^T=-Sigma_W
% into an ordinary linear system (I(x)A+A(x)I)vec(Sigma)=-vec(Sigma_W), then
% show that Kronecker-sum matrix is invertible (Hurwitz => nonsingular).

Re-arranging the terms in equation \ref{eq:lyapunov_eq0} and let $\mathbf{\Sigma}_{\mathbf{X}(\infty)}=\mathbf{\Sigma}_{\infty}$, we get
\begin{equation}
\mathbf{A} \mathbf{\Sigma}_\mathbf{\infty} +\mathbf{\Sigma}_\mathbf{\infty} \mathbf{A}^T=-\mathbf{\Sigma}_\mathbf{W}
\label{eq:ss_lya_re0}
\end{equation}


Applying the vectorization operator $\operatorname{vec}$, which stacks the columns of a matrix to form a vector \citep{MagnusNeudecker1999}, to equation \ref{eq:ss_lya_re0}, we get.
\begin{equation}
\operatorname{vec}(\mathbf{A} \mathbf{\Sigma}_\mathbf{\infty} +\mathbf{\Sigma}_\mathbf{\infty} \mathbf{A}^T)=-\operatorname{vec} \mathbf{\Sigma}_\mathbf{W}
\label{eq:ss_ly_re1}
\end{equation}


Or equivalently,
\begin{equation}
\operatorname{vec}(\mathbf{A} \mathbf{\Sigma}_\mathbf{\infty} \mathbf{I})+\operatorname{vec}(\mathbf{I} \mathbf{\Sigma}_\mathbf{\infty} \mathbf{A}^T)=
-\operatorname{vec} \mathbf{\Sigma}_\mathbf{W}
\label{eq:ss_ly_re2}
\end{equation}


The following identity holds for matrices $A $, $B $, and $C $ that are conformal with respect to each other, where $\otimes$ is the Kronecker product \citep{MagnusNeudecker1999}.
\begin{equation}
\operatorname{vec}(\mathbf{ABC})=(\mathbf{C}^T \otimes \mathbf{A} )\operatorname{vec}(\mathbf{B})
\label{eq:vec_formula}
\end{equation}


Applying the formula from equation~\ref{eq:vec_formula} to equation~\ref{eq:ss_ly_re2}, we get
\begin{equation}
(\mathbf{I}^T \otimes \mathbf{A}) \operatorname{vec} \mathbf{\Sigma}_\mathbf{\infty} + \left ((\mathbf{A}^T)^T \otimes \mathbf{I} \right) \operatorname{vec} \mathbf{\Sigma}_\mathbf{\infty} = -\operatorname{vec} \mathbf{\Sigma}_\mathbf{W}
\end{equation}


Or equivalently,
\begin{equation}
(\mathbf{I} \otimes \mathbf{A} + \mathbf{A} \otimes \mathbf{I}) \operatorname{vec} \mathbf{\Sigma}_\mathbf{\infty} = -\operatorname{vec} \mathbf{\Sigma}_\mathbf{W}
\end{equation}


% Nonsingularity argument: Hurwitz A => all pairwise eigenvalue sums have
% negative real part => Kronecker sum has no zero eigenvalue => invertible.
The matrix $\mathbf{I} \otimes \mathbf{A} + \mathbf{A} \otimes \mathbf{I}$ is the Kronecker sum $\mathbf{A} \oplus \mathbf{A}$, whose eigenvalues are exactly the pairwise sums $\lambda_i + \lambda_j$ of the eigenvalues $\{\lambda_i\}$ of $\mathbf{A}$ \citep[Theorem~4.4.5]{HornJohnson1991}. Because $\mathbf{A}$ is Hurwitz, $\mathrm{Re}(\lambda_i) < 0$ for all $i$, so $\mathrm{Re}(\lambda_i + \lambda_j) < 0$ for every pair $(i,j)$; in particular no eigenvalue is zero, and the Kronecker sum is therefore nonsingular. Hence it is invertible, and
\begin{equation}
\operatorname{vec} \mathbf{\Sigma}_\mathbf{\infty} =
-{(\mathbf{I} \otimes \mathbf{A} + \mathbf{A} \otimes \mathbf{I})}^{-1} \operatorname{vec}(\mathbf{\Sigma}_\mathbf{W})
\label{eq:sigma_infty_vec}
\end{equation}


\subsection*{An Analytic Relationship between Entries of $\mathbf{\Sigma_Y}$ and $A$}
% Substitute the vec(Sigma_infty) solution (eq:sigma_infty_vec) into
% vec(Sigma_Y) to express the observed covariance directly in terms of A
% and Sigma_W.
Applying the vectorization operator to equation \ref{eq:cov_yt}, we get.
\begin{equation}
\begin{aligned}
\operatorname{vec} \mathbf{\Sigma_Y} &= \operatorname{vec} (\mathbf{H} \mathbf{\Sigma}_\mathbf{\infty} \mathbf{H}^T) + \operatorname{vec} \mathbf{\Sigma}_V =\left( (\mathbf{H}^T)^T \otimes \mathbf{H} \right)\operatorname{vec} \mathbf{\Sigma}_\mathbf{\infty} + \operatorname{vec} \mathbf{\Sigma}_V =\left(\mathbf{H} \otimes \mathbf{H} \right)\operatorname{vec} \mathbf{\Sigma}_\mathbf{\infty} + \operatorname{vec} \mathbf{\Sigma}_V
\end{aligned}
\label{eq:cov_yt_vec}
\end{equation}


Substituting $\operatorname{vec} \mathbf{\Sigma}_\infty$ from equation \ref{eq:sigma_infty_vec} into equation \ref{eq:cov_yt_vec}, we get.
\begin{equation}
\begin{aligned}
\operatorname{vec} \mathbf{\Sigma_Y} =\left(\mathbf{H} \otimes \mathbf{H} \right) \operatorname{vec} \mathbf{\Sigma}_\mathbf{\infty} + \operatorname{vec} \mathbf{\Sigma}_V =-\left(\mathbf{H} \otimes \mathbf{H} \right) \left \{ {(\mathbf{I} \otimes \mathbf{A} + \mathbf{A} \otimes \mathbf{I})}^{-1} \operatorname{vec}\mathbf{\Sigma}_\mathbf{W} \right \} + \operatorname{vec} \mathbf{\Sigma}_V
\end{aligned}
\label{eq:cov_yt_A_vec}
\end{equation}


% Specializes to the fully-observed, noise-free case used throughout the
% rest of the paper (Sigma_Y is then just the state covariance Sigma_infty).
With very little loss of generality, we will assume $\mathbf{H}=\mathbf{I}$ and $\mathbf{\Sigma}_V=\mathbf{0}$, i.e., the states $\mathbf{X}(t)$ of the dynamical system are directly observable and there is no observation noise $\mathbf{V} (t)$. The equation above (\ref{eq:cov_yt_A_vec}) can then be simplified as
\begin{equation}
\begin{aligned}
\operatorname{vec} \mathbf{\Sigma_Y} &=- {(\mathbf{I} \otimes \mathbf{A} + \mathbf{A} \otimes \mathbf{I})}^{-1} \operatorname{vec} \mathbf{\Sigma}_\mathbf{W}
\end{aligned}
\label{eq:cov_yt_A_vec1}
\end{equation}


% From here on, work out vec(Sigma_Y)=-K^{-1}vec(Sigma_W) (eq:cov_yt_A_vec1)
% explicitly for a general n-node network, to show sigma_ij/rho_ij are
% closed-form (rational) functions of A's entries -- not yet solving any
% specific topology (that's Appendices C/D).
% describe a N x N network
Let us consider an $N \times N$ network with the following architecture.
% N by N matrix
% a slightly more compact version of these matrices
\begin{equation}
\begin{aligned}
\mathbf{A} =
\begin{bmatrix}
    a_{11}  & \cdots & a_{1n} \\
     \vdots & \ddots & \vdots  \\
    a_{n1}  & \cdots & a_{nn}
\end{bmatrix},
\mathbf{\Sigma_\mathbf{Y}}=
\begin{bmatrix}
    \sigma_{11}  & \cdots & \sigma_{1n} \\
     \vdots & \ddots & \vdots  \\
    \sigma_{n1}  & \cdots & \sigma_{nn}
\end{bmatrix},
\mathbf{\Sigma_\mathbf{W}} =
\begin{bmatrix}
    \sigma_{W_1}  & \cdots & 0 \\
     \vdots & \ddots & \vdots  \\
    0  & \cdots & \sigma_{W_n}
\end{bmatrix}
\end{aligned}
\label{eq:nByn_network_append}
\end{equation}


% This is the same Kronecker-sum matrix from eq:sigma_infty_vec, just
% written out explicitly in n-by-n block form for this general network.
Therefore,
\begin{equation}
\begin{aligned}
\mathbf{I} \otimes \mathbf{A}+ \mathbf{A} \otimes \mathbf{I} &=
\left[ \begin{array}{cccc}
\mathbf{A} & \mathbf{0} & \cdots & \mathbf{0} \\
\mathbf{0} & \mathbf{A} & \cdots & \mathbf{0} \\
\vdots & \vdots & \ddots & \vdots \\
\mathbf{0} & \mathbf{0} & \cdots & \mathbf{A}
\end{array} \right] +
\left[ \begin{array}{cccc}
a_{11}\mathbf{I} & a_{12}\mathbf{I} & \cdots & a_{1n}\mathbf{I} \\
a_{21}\mathbf{I} & a_{22}\mathbf{I} & \cdots & a_{2n}\mathbf{I} \\
\vdots & \vdots & \ddots & \vdots \\
a_{n1}\mathbf{I} & a_{n2}\mathbf{I} & \cdots & a_{nn}\mathbf{I}
\end{array} \right] \\
&=\left[ \begin{array}{cccc}
\mathbf{A}+a_{11}\mathbf{I} & a_{12}\mathbf{I} & \cdots & a_{1n}\mathbf{I} \\
a_{21}\mathbf{I} & \mathbf{A}+a_{22}\mathbf{I} & \cdots & a_{2n}\mathbf{I} \\
\vdots & \vdots & \ddots & \vdots \\
a_{n1}\mathbf{I} & a_{n2}\mathbf{I} & \cdots & \mathbf{A}+a_{nn}\mathbf{I}
\end{array} \right] \\
&=\left[ \begin{array}{cccc}
k_{11}^{11} & k_{21}^{11} & \cdots & k_{nn}^{11} \\
k_{11}^{21} & k_{21}^{21} & \cdots & k_{nn}^{21} \\
\vdots & k_{ij}^{uv} & \ddots & \vdots \\
k_{11}^{nn} & k_{21}^{nn} & \cdots & k_{nn}^{nn} \\
\end{array} \right]=\mathbf{K}_{n^2 \times n^2}
\end{aligned}
\end{equation}



% Key existence argument: |K| is a polynomial in K's own entries (Leibniz
% formula), and each entry of K is itself linear in A's entries, so |K| is
% a well-defined polynomial function g of A's entries too -- this is what
% guarantees a genuine closed form exists, before deriving its details.
Where we use superscripts to indicate the row indices and subscripts the column indices of $\mathbf{K}$, respectively. According to the Leibniz formula, the determinant of $\mathbf{K}$ is a homogeneous polynomial of degree $n^2$ in the entries of $\mathbf{K}$: each of the $(n^2)!$ signed terms is a product of $n^2$ entries, one taken from each row and each column. Hence, it is a well-defined (polynomial, and therefore nonlinear) function, which we represent with $g_{0}$

\begin{equation}
|\mathbf{K}|=g_{0}(k_{11}^{11},k_{21}^{11}, k_{ij}^{uv},\cdots, k_{nn}^{nn})
\label{eq:det_k}
\end{equation}

% John P. thinks the following explanation is not necessary. I think so too.
\begin{comment}
\sum_{\sigma \in S_{n^2}} sgn(\sigma) k_{\sigma(11)}^{11} k_{\sigma(21)}^{21}, \cdots, k_{\sigma(nn)}^{nn}

Where $\sigma$ is a bijective function from the set of column indices ${11, 21, \cdots, nn }$ to itself, i.e., the output of the function is a permutation of the set, and $sgn(\sigma)$ is the signature of the permutation. The signature $+1$ is  if the permutation can be obtained via an even number of transpositions (swapping two elements of the set). Otherwise, the signature is $-1$,  if the permutation can only be obtained via an odd number of swaps. Finally, $S_{n^2}$ is the symmetry group of the set ${11, 21, \cdots, nn }$, which contains all such permutations. The important thing to recognize here is that the determinant of a square matrix is always well defined and it is the sum of the signed products of entries chosen from the matrix, where each entry is chosen from a unique row such that no two entries are from the same column. Thus, each summand is a polynomial function of n entries from the matrix $K$. Therefore, the determinant of $K$ can be considered as a nonlinear function of its entries and we write it as the following,

\begin{equation}
|K|=\sum_{\sigma \in S_{n^2}} sgn(\sigma) k_{\sigma(11)}^{11} k_{\sigma(21)}^{21}, \cdots, k_{\sigma(nn)}^{nn}=g_{0}(k_{11}^{11},k_{21}^{11}, k_{ij}^{uv},\cdots, k_{nn}^{nn})
\end{equation}
\end{comment}


Because each entry ($k_{ij}^{uv}$) of the matrix $\mathbf{K}=\mathbf{I} \otimes \mathbf{A}+ \mathbf{A} \otimes \mathbf{I}$ is a linear combination of the entries of $\mathbf{A}$ (a single off-diagonal entry $a_{ij}$, a sum $a_{ii}+a_{kk}$ on a diagonal position, or $0$), composing $g_{0}$ with these linear maps shows that $|\mathbf{K}|$ is also a polynomial function, and hence a nonlinear function ($g$), of the entries of $\mathbf{A}$.
\begin{equation}
|\mathbf{K}|=
g_{0}(k_{11}^{11},k_{21}^{11}, k_{ij}^{uv},\cdots, k_{nn}^{nn})=g(a_{11},a_{12},\cdots, a_{nn})
\end{equation}


The inverse of $\mathbf{K}$ can be computed from its cofactor matrix $\mathbf{C}$ and determinant $|\mathbf{K}|$.
\begin{equation}
\mathbf{K}^{-1}=\frac{1}{|\mathbf{K}|} \mathbf{C}^T=\frac{1}{|\mathbf{K}|}
\left[ \begin{array}{cccc}
C_{11}^{11} & C_{21}^{11} & \cdots & C_{nn}^{11} \\
C_{11}^{21} & C_{21}^{21} & \cdots & C_{nn}^{21} \\
\vdots & C_{ij}^{uv} & \ddots & \vdots \\
C_{11}^{nn} & C_{21}^{nn} & \cdots & C_{nn}^{nn} \\
\end{array} \right]^T
\end{equation}


\begin{comment}
\frac{1}{|\mathbf{K}|}\left[ \begin{array}{cccc}
C_{11}^{11} & C_{11}^{21} & \cdots & C_{11}^{nn} \\
C_{21}^{11} & C_{21}^{21} & C_{ij}^{uv} & C_{21}^{nn} \\
\vdots & \vdots & \ddots & \vdots \\
C_{nn}^{11} & C_{nn}^{21} & \cdots & C_{nn}^{nn} \\
\end{array} \right]
\end{comment}


Each entry $C_{ij}^{uv}$ of matrix $\mathbf{C}$ is obtained by multiplying $(-1)^{\{u+(v-1)n\}+\{i+(j-1)n\}}$ by the associated minor $M_{\{u+(v-1)n\},\,\{i+(j-1)n\}}$ of $\mathbf{K}$, i.e.
\begin{comment}
\begin{equation}
C=\left[ \begin{array}{cccc}
C_{11}^{11} & C_{21}^{11} & \cdots & C_{nn}^{11} \\
C_{11}^{21} & C_{21}^{21} & \cdots & C_{nn}^{21} \\
\vdots & \vdots & \ddots & \vdots \\
C_{11}^{nn} & C_{21}^{nn} & \cdots & C_{nn}^{nn} \\
\end{array} \right]
\end{equation}
\end{comment}

\begin{equation}
C_{ij}^{uv}=(-1)^{\{u+(v-1)n\}+\{i+(j-1)n\}}M_{\{u+(v-1)n\},\,\{i+(j-1)n\}}
\end{equation}


Where $i=1,2,\cdots,n$, $j=1,2,\cdots,n$, $u=1,2,\cdots,n$, and $v=1,2,\cdots, n$. Each pair of $(i,j)$ identifies a unique column of the cofactor matrix $\mathbf{C}$  and each pair of $(u,v)$ a unique row of $\mathbf{C}$, respectively. The minor $M_{\{u+(v-1)n\},\,\{i+(j-1)n\}}$ is the determinant of the submatrix formed by deleting the $\{u+(v-1)n\}$ row and $\{i+(j-1)n\}$ column of matrix $\mathbf{K}$. For example,
\begin{equation}
\left[ \begin{array}{ccc}
a & d & g \\
b & e & h \\
c & f & i \\
\end{array} \right],
\quad
M_{2,3}=det
\left[ \begin{array}{ccc}
a & d & {[]} \\
{[]} & {[]} & {[]} \\
c & f & {[]} \\
\end{array} \right]=
det
\left[ \begin{array}{cc}
a & d \\
c & f \\
\end{array} \right]=af-cd
\end{equation}


Again, each entry of the cofactor matrix $\mathbf{C}$ is well defined and can be identified as a nonlinear function of $(k_{11}^{11},k_{21}^{11}, k_{ij}^{uv},\cdots, k_{nn}^{nn})$ or equivalently a nonlinear function of $(a_{11},a_{12},\cdots, a_{nn})$.
%%%%%%%%%%%%%%%%%%%%%
\begin{equation}
C_{ij}^{uv}=f_{ij}^{uv}(a_{11},a_{12},\cdots, a_{nn})
\label{eq:cofactor_poly}
\end{equation}
%%%%%%%%%%%%%%%%%%%%%

Equation $\operatorname{vec} \mathbf{\Sigma}_Y =- {(\mathbf{I} \otimes \mathbf{A} + \mathbf{A} \otimes \mathbf{I})}^{-1} \operatorname{vec} \mathbf{\Sigma}_\mathbf{W}$ (\ref{eq:cov_yt_A_vec1}) can now be written, equivalently, as the following.

\begin{equation}
\operatorname{vec} \mathbf{\Sigma}_Y=
\left[ \begin{array}{c}
\sigma_{11} \\
\vdots \\
\sigma_{n1} \\
\sigma_{ij} \\
\vdots \\
\sigma_{n2} \\
\sigma_{1n} \\
\vdots \\
\sigma_{nn} \\
\end{array} \right]=
-\mathbf{K}^{-1} \operatorname{vec} \mathbf{\Sigma_W} =- \frac{1}{|\mathbf{K}|}
\left[ \begin{array}{cccc}
C_{11}^{11} & C_{21}^{11} & \cdots & C_{nn}^{11} \\
C_{11}^{21} & C_{21}^{21} & \cdots & C_{nn}^{21} \\
\vdots & C_{ij}^{uv} & \ddots & \vdots \\
C_{11}^{nn} & C_{21}^{nn} & \cdots & C_{nn}^{nn} \\
\end{array} \right]^T
\left[ \begin{array}{c}
\sigma_{W_1} \\
0 \\
\vdots \\
0 \\
\sigma_{W_2} \\
0 \\
\vdots \\
0 \\
\sigma_{W_n} \\
\end{array} \right]
\end{equation}


Or equivalently,
\begin{equation}
\sigma_{ij}=\sigma_{ji}=- \frac{1}{|\mathbf{K}|}
\sum_{u=1}^{n} \sum_{v=1}^{n} C_{ij}^{uv} \sigma_{W_{uv}}=
- \frac{\sum_{u=1}^{n} \sum_{v=1}^{n}
f_{ij}^{uv}(a_{11},a_{12},\cdots, a_{nn}) \sigma_{W_{uv}}}{g(a_{11},a_{12},\cdots, a_{nn})}
\label{eq:sigma_ij_final}
\end{equation}

Each pair $(i,j)$ identifies a unique entry of the covariance matrix $\mathbf{\Sigma_Y}$, which corresponds to a unique column of the cofactor matrix $\mathbf{C}$, and each pair $(u,v)$ identifies a unique row of $\mathbf{C}$. Summing over $u$ and $v$ therefore cycles through all the entries (rows) in the column identified by $(i,j)$. Because $\mathbf{\Sigma_W}$ is diagonal, $\sigma_{W_{uv}}=0$ whenever $u \neq v$, so only the diagonal terms ($u=v$) contribute to these sums. The correlation coefficient then follows from its definition, $\rho_{ij}=\sigma_{ij}/\sqrt{\sigma_{ii}\,\sigma_{jj}}$. 
\begin{equation}
\begin{aligned}
\rho_{ij} &=\frac{\sigma_{ij}}{\sqrt{\sigma_{ii}\,\sigma_{jj}}} \\
\\
&=\frac{- \frac{1}{|\mathbf{K}|}\sum_{u=1}^{n} C_{ij}^{uu} \sigma_{W_{u}}}{\sqrt{\left(- \frac{1}{|\mathbf{K}|}\sum_{u=1}^{n} C_{ii}^{uu} \sigma_{W_{u}}\right)\left(- \frac{1}{|\mathbf{K}|}\sum_{u=1}^{n} C_{jj}^{uu} \sigma_{W_{u}}\right)}} \\
\\
&=-\operatorname{sgn}(\det \mathbf{K})\, \frac{ \sum_{u=1}^{n} C_{ij}^{uu} \sigma_{W_{u}}}{\sqrt{\left( \sum_{u=1}^{n} C_{ii}^{uu} \sigma_{W_{u}}\right)\left( \sum_{u=1}^{n} C_{jj}^{uu} \sigma_{W_{u}}\right)}}
\end{aligned}
\label{eq:rho_ij_nbyn}
\end{equation}

where the cofactor terms $C_{ij}^{uu}$ are the polynomials in the causal parameters given by equation~(\ref{eq:cofactor_poly}).
% This is the appendix's central result: rho_ij, for an arbitrary n-node
% network, is a ratio of polynomials in A's entries and the noise
% variances -- the general existence claim that Appendices C/D make
% concrete for n=2 and n=3.


%%%%%%%%%%%%



\bibliographystyle{apacite}
\bibliography{appendix_b}

\end{document}